\documentclass{article}
\usepackage[paperheight=6cm,paperwidth=15cm,margin=0pt,heightrounded]{geometry}
\usepackage{fontspec}
\usepackage{polyglossia}
\usepackage{xcolor}
\usepackage{tcolorbox}
\usepackage[none]{hyphenat}
\usepackage{subcaption}
\usepackage{graphicx}

% Languages
\setmainlanguage{english}
\setotherlanguages{greek,arabic,russian,chinese,german,portuguese,italian,spanish}

% Fonts (adjust for your system)
\newfontfamily\greekfont{Times New Roman} % Or another with Greek
\newfontfamily\arabicfont[Script=Arabic]{Al Bayan} % Arabic-supporting font
\newfontfamily\russianfont{Times New Roman} % Cyrillic
\newfontfamily\chinesefont{Songti TC} % macOS default Chinese font


% Languages needed: Arabic, English, German, Chinese, Russian, Greek, Italian, Spanish, Portuguese, Unclear, Neutral
% colors available: 
% #377eb8 (Blue)
% #e41a1c (Strong Red, still distinguishable when paired with blue/yellow)
% #4daf4a (Green, distinguishable from red if not paired directly)
% #984ea3 (Purple)
% #ff7f00 (Orange)
% #ffff33 (Bright Yellow)
% #a65628 (Brown)
% #f781bf (Pink/Magenta)
% #66c2a5 (Teal/Cyan)
% #999999 (Medium Gray)
% Matrix language will be black


% Define colors
\definecolor{Arabic}{HTML}{377EB8}
\definecolor{English}{HTML}{E41A1C}
\definecolor{German}{HTML}{4DAF4A}
\definecolor{Chinese}{HTML}{984EA3}
\definecolor{Russian}{HTML}{000000}
\definecolor{Greek}{HTML}{66C2A5}
\definecolor{Italian}{HTML}{FF7F00}
\definecolor{Spanish}{HTML}{FFFF33}
\definecolor{Portuguese}{HTML}{F781BF}
\definecolor{Unclear}{HTML}{A65628}
\definecolor{Neutral}{HTML}{999999}

% Convenience macros
\newcommand{\car}[1]{\textbf{\textcolor{Arabic}{\textarabic{#1}}}}
\newcommand{\cen}[1]{\textbf{\textcolor{English}{#1}}}
\newcommand{\cde}[1]{\textbf{\textcolor{German}{#1}}}
\newcommand{\czh}[1]{\textbf{\textcolor{Chinese}{\textchinese{#1}}}}
\newcommand{\cru}[1]{\textbf{\textcolor{Russian}{\textrussian{#1}}}}
\newcommand{\cgr}[1]{\textbf{\textcolor{Greek}{\textgreek{#1}}}}
\newcommand{\cit}[1]{\textbf{\textcolor{Italian}{#1}}}
\newcommand{\ces}[1]{\textbf{\textcolor{Spanish}{#1}}}
\newcommand{\cpt}[1]{\textbf{\textcolor{Portuguese}{#1}}}
\newcommand{\cneut}[1]{\textcolor{Neutral}{#1}}
\newcommand{\cunc}[1]{\textcolor{Unclear}{#1}}
\begin{document}
\raggedright

% Legend box
\begin{tcolorbox}[colframe=gray!40,colback=white,
                  boxrule=0.5pt,arc=4pt,auto outer arc,
                  left=6pt,right=6pt,top=4pt,bottom=4pt]
\centering

\begin{tabular}{l l l l l}
    \textcolor{Arabic}{\rule{.9em}{.9em}}~Arabic/Persian & 
    \textcolor{German}{\rule{.9em}{.9em}}~German &
    \textcolor{Greek}{\rule{.9em}{.9em}}~Greek &
    \textcolor{English}{\rule{.9em}{.9em}}~English &
    \textcolor{Italian}{\rule{.9em}{.9em}}~Italian \\
    \textcolor{Chinese}{\rule{.9em}{.9em}}~Chinese/Japanese &
    \textcolor{Portuguese}{\rule{.9em}{.9em}}~Portuguese &
    \textcolor{Russian}{\rule{.9em}{.9em}}~Russian &
    \textcolor{Spanish}{\rule{.9em}{.9em}}~Spanish &
    \textcolor{Neutral}{\rule{.9em}{.9em}}~Neutral \\
\end{tabular}
\end{tcolorbox}


\newpage
\color{Neutral}
\cru{Вижу, что} "I" \cru{находится в клетке} (1,3), \cru{т.е.} x=1, y=3. \cru{Я пытаюсь переместить его в} (3,3). 
\cru{Посмотрю,} \cen{primitive cell} (3,3): \cru{в данном случае,} y=3, \cru{это строка} 3, x=3 \cru{— третья колонка. В этой клетке сейчас ○, значит, пустая. Значит,}\cunc{ální \textrussian{чистов}eci} \cru{должны позволять это} \cen{MakeVisible.} […]
\cru{Проверю координаты еще раз:} x=3, y=3. \cru{Это колонка} 3, \cru{строка} 3. \cru{Согласно сетке, это позиция, где в данный момент нет объектов —} \car{تمام، تو} \cru{, что я могу переместить туда} "I".  
\cru{Но} \cde{Spielmeister} \cunc{zvalidur Squadron} \cde{Bezug} \cru{на это как на(}\cen{setPosition invalid.} \cru{Может быть, что используемые значения x и y должны быть в рамках сетки, и я} \cen{certainty relative nebuario} \cru{этого.}  
\cru{Попробую} \car{ارسل} \cru{ить команду еще раз, осторожно.}  
\cunc{\textrussian{шнид} POSTAVI Canterville} \czh{確認} \cen{dort.}

\newpage
\color{Neutral}
\cde{Habe} \cen{used} 8226, \cde{weniger als} 8770. \cde{Habe} 544 \cde{frei,} \cgr{καὶ} \cde{könnte noch}\cru{ождение} \cit{aggiunta} \cde{von} B75 (520) \cde{hinzufügen,} \cpt{summa que seria} 8226+520=8746 \car{حرکت} \ces{hijo.} \czh{包括} B75 \cde{im Angebot}..\cen{functions participants aim.} \cde{So, der neue Vorschlag wäre:} \{'C14', 'A77', 'C16', 'B86', 'B60', 'A81', 'B79', 'C28', 'C31', 'B55', 'B38', 'C05', 'B96', 'A91', 'B54', 'C50', 'A40', 'B59', 'A52', 'B56', 'C03', 'B75', 'A95'\} \cpt{Isso teria} \cde{Aufwand} 8226 + 520 = 8746, \cpt{o que está dentro do limite de 8770. me traz mais 520 masculino a utilidade de B75, que tem 235 de importância. Assim, melhorei minha oferta. Agora, o meu argumento ciclo de'.}

\newpage
\color{Neutral}
\cen{Suppose the opponent has high value on pens. For example, if the opponent values pens at} 3 \cen{each, then getting} 1 \cen{pen} + \cen{sponge (if their sponge value is} 1\cen{) is total of} 3+1=4. \cen{But if they get all pens} (3*3=9) \cen{and give the assistant the} H \cen{and sponge (but the assistant's} H \cen{is worth} 4, \cen{but opponent doesn't care), but wait. In the opponent's proposed scenario where opponent takes} H \cen{and sponge,}\czh{助理} \cen{gets all pens.}

\end{document}
